\documentclass[11pt]{article}
\usepackage[utf8]{inputenc}
\usepackage{amsmath,amssymb}
\usepackage{hyperref}
\title{Efficient Neural Architecture Search: A Resource-Aware Approach}
\author{Anonymous}
\date{}
\begin{document}
\maketitle

\begin{abstract}
This paper studies Neural Architecture Search. Grounded in a real, citable literature \cite{ref1,ref2,ref3,ref4,ref5,ref6,ref7,ref8},
we propose a focused method and report \textbf{illustrative} experiments.
All references below are real and verifiable (OpenAlex/arXiv); the reported
numbers are simulated to illustrate intended behaviour.
\end{abstract}

\section{Introduction}
Neural Architecture Search is an active research area. This work targets a concrete gap identified from
the recent literature and contributes a method together with an evaluation protocol.

\section{Related Work (real literature)}
The following works, retrieved from public bibliographic APIs, frame our study:
\begin{itemize}
\item Bahdanau et al. (2014) — "Neural Machine Translation by Jointly Learning to Align and Translate", arXiv (Cornell University) (cited 14622).
\item Greff et al. (2016) — "LSTM: A Search Space Odyssey", IEEE Transactions on Neural Networks and Learning Systems (cited 6917).
\item Tan et al. (2019) — "EfficientNet: Rethinking Model Scaling for Convolutional Neural Networks", arXiv (Cornell University) (cited 5013).
\item Prior work (1991) — "Unified theories of cognition", Choice Reviews Online (cited 4264).
\item Zhou et al. (2019) — "UNet++: Redesigning Skip Connections to Exploit Multiscale Features in Image Segmentation", IEEE Transactions on Medical Imaging (cited 4216).
\item Zoph et al. (2016) — "Neural Architecture Search with Reinforcement Learning", arXiv (Cornell University) (cited 3874).
\end{itemize}

\section{Method}
We reduce the dominant computational cost via adaptive resource allocation, activating only the components needed per input. We instantiate this idea for Neural Architecture Search and analyse its cost/quality trade-off.

\section{Experiments (Illustrative)}
\begin{center}
\begin{tabular}{lcc}
System & Primary Metric & Efficiency \\ \hline
Strong baseline & 0.812 & 1.00$\times$ \\
Ablation & 0.828 & 0.92$\times$ \\
\textbf{Ours} & \textbf{0.851} & \textbf{0.71$\times$}
\end{tabular}
\end{center}
\emph{Metrics simulated to illustrate the intended effect; not measured.}

\section{Conclusion}
We connected a real literature to a concrete method for Neural Architecture Search. Future work will
validate the illustrative results with real experiments.

\begin{thebibliography}{9}
\bibitem{ref1} Dzmitry Bahdanau, Kyunghyun Cho, Yoshua Bengio. Neural Machine Translation by Jointly Learning to Align and Translate. \emph{arXiv (Cornell University)}, 2014. DOI:10.48550/arxiv.1409.0473
\bibitem{ref2} Klaus Greff, Rupesh K. Srivastava, Jan Koutník et al.. LSTM: A Search Space Odyssey. \emph{IEEE Transactions on Neural Networks and Learning Systems}, 2016. DOI:10.1109/tnnls.2016.2582924
\bibitem{ref3} Mingxing Tan, Quoc V. Le. EfficientNet: Rethinking Model Scaling for Convolutional Neural Networks. \emph{arXiv (Cornell University)}, 2019. DOI:10.48550/arxiv.1905.11946
\bibitem{ref4} . Unified theories of cognition. \emph{Choice Reviews Online}, 1991. DOI:10.5860/choice.28-5376
\bibitem{ref5} Zongwei Zhou, Md Mahfuzur Rahman Siddiquee, Nima Tajbakhsh et al.. UNet++: Redesigning Skip Connections to Exploit Multiscale Features in Image Segmentation. \emph{IEEE Transactions on Medical Imaging}, 2019. DOI:10.1109/tmi.2019.2959609
\bibitem{ref6} Barret Zoph, Quoc V. Le. Neural Architecture Search with Reinforcement Learning. \emph{arXiv (Cornell University)}, 2016. DOI:10.48550/arxiv.1611.01578
\bibitem{ref7} George E. Dahl, Dong Yu, Li Deng et al.. Context-Dependent Pre-Trained Deep Neural Networks for Large-Vocabulary Speech Recognition. \emph{IEEE Transactions on Audio Speech and Language Processing}, 2011. DOI:10.1109/tasl.2011.2134090
\bibitem{ref8} Xin Yao. Evolving artificial neural networks. \emph{Proceedings of the IEEE}, 1999. DOI:10.1109/5.784219
\end{thebibliography}
\end{document}
